\usepackage{amsmath}
\usepackage{amssymb}
\usepackage{multirow}
\usepackage{chemfig}
\usetikzlibrary{arrows.meta,positioning,shapes.geometric,patterns,calc,decorations.pathreplacing}

% \qui{SiO_2} escreve a fórmula com o índice subscrito.
% Feito à mão para não depender do mhchem.
\newcommand{\qui}[1]{\ensuremath{\mathrm{#1}}}
\newcommand{\seta}{\ensuremath{\;\longrightarrow\;}}
\newcommand{\setaeq}{\ensuremath{\;\rightleftharpoons\;}}

% No lugar do siunitx.
\newcommand{\un}[2]{\ensuremath{#1\,\mathrm{#2}}}
\newcommand{\grau}[1]{\ensuremath{#1\,^{\circ}\mathrm{C}}}

\newcommand{\termo}[1]{\textcolor{mPrimaryDark}{\textbf{#1}}}

% Caixa que fecha um tópico ligando o assunto a telecom.
\newcommand{\gancho}[1]{%
  \begin{alertblock}{Por que isso importa em Telecomunicações?}#1\end{alertblock}}

\newcommand{\fonte}[1]{%
  \vfill{\scriptsize\textcolor{mNeutral}{Fonte: #1}}}

% \atomo{x}{y}{cor}{rótulo}: átomo dos diagramas de rede.
\newcommand{\atomo}[4]{%
  \node[circle,draw=#3,fill=#3!20,thick,minimum size=0.95cm,inner sep=0pt,
        font=\scriptsize\bfseries,text=#3] at (#1,#2) {#4};}
